% Conexões do Subcomponente 3.2.b com outras partes do MOP
\begin{tikzpicture}[
    caixa/.style={draw=#1, fill=#1!8, rounded corners=3pt, align=left,
                  text width=4.1cm, inner sep=5pt, font=\scriptsize,
                  minimum height=1.55cm},
    entra/.style={-{Stealth[length=2.4mm]}, line width=1.1pt, draw=#1},
    sai/.style={{Stealth[length=2.4mm]}-, line width=1.1pt, draw=#1}
  ]
  \node[marco=3.6cm, minimum height=2.2cm, font=\footnotesize\bfseries] (c) at (0,0)
    {Subcomponente 3.2.b\\[2pt]\normalfont\scriptsize abastecimento e esgotamento no meio rural};

  % entradas (à esquerda)
  \node[caixa=azulpsh] (a) at (-5.6,2.6) {\textbf{3.2.a -- Política e Plano Estadual de
    Saneamento Rural} (SEPL/SECID)\\[2pt]\textit{fornece diretrizes, modelos de gestão e
    critérios de priorização}};
  \node[caixa=green!45!black] (b) at (-5.6,0) {\textbf{2.3 -- Agricultura sustentável em
    bacias de mananciais} (SEAB/IDR-Paraná)\\[2pt]\textit{os PIP indicam os domicílios da
    frente mananciais}};
  \node[caixa=teal!80!black] (d) at (-5.6,-2.6) {\textbf{1.1 -- Instrumentos de gestão e
    monitoramento} (IAT)\\[2pt]\textit{outorga do uso da água dos poços e registro de
    dados no SIGARH}};

  % saídas e apoio (à direita)
  \node[caixa=cyan!60!black] (e) at (5.6,2.6) {\textbf{3.1 -- Gestão de riscos em
    mananciais} (Sanepar)\\[2pt]\textit{recebe menos esgoto sem tratamento nas bacias de
    mananciais}};
  \node[caixa=orange!85!black] (f) at (5.6,0) {\textbf{Indicadores do Componente 3}\\[2pt]
    \textit{recebe pessoas atendidas com água e esgoto e associações com modelo de gestão}};
  \node[caixa=gray] (g) at (5.6,-2.6) {\textbf{Componente 4 -- Gestão do projeto}
    (UGP/SEPL)\\[2pt]\textit{apoia aquisições, desembolsos, monitoramento e reporte ao
    Banco Mundial}};

  \draw[entra=azulpsh] (a.east) to[out=0, in=150] (c.west);
  \draw[entra=green!45!black] (b.east) -- (c.west);
  \draw[entra=teal!80!black] (d.east) to[out=0, in=210] (c.west);
  \draw[entra=cyan!60!black] (c.east) to[out=30, in=180] (e.west);
  \draw[entra=orange!85!black] (c.east) -- (f.west);
  \draw[entra=gray, <->] (c.east) to[out=330, in=180] (g.west);

  \node[font=\scriptsize\bfseries, text=azulpsh] at (-5.6,4.0) {O que o 3.2.b recebe};
  \node[font=\scriptsize\bfseries, text=azulpsh] at (5.6,4.0) {O que o 3.2.b entrega ou troca};
\end{tikzpicture}
